\begin{lemma}[Old Algorithm 3.1 steps survive a data append]
Fix finite graph data, a color bound \(k\), distinguished edge \(e_0\), input
coloring \(\varphi\), and Tashkinov-series data \(A\) with \(r\) noninitial
iterations.  Append arbitrary supplied final values to obtain \(A^+\) as in
\(\noderef{TashkinovSeriesDataAppend}\).  If an old index
\(j\in\operatorname{Fin}(r+1)\) has an Algorithm 3.1 step certificate for
\((G,k,r,e_0,\varphi,A)\), then its canonical image \(j^+\) has an Algorithm
3.1 step certificate for \((G,k,r+1,e_0,\varphi,A^+)\).
\end{lemma}
\begin{proof}
Use the exact old-slot equalities from
\(\noderef{TashkinovSeriesDataAppendEquations}\).  Besides their stated
cast-successor form, the elementary index equalities
\((i^+).succ=(i.succ)^+\) and their tree-index analogues identify every tree
and non-tree expression occurring in an old constructor with the expression
at the lifted constructor.  We inspect the four constructors of
\(\noderef{TashkinovSeriesStep}\).

For the initial constructor, the new indices zero and one are the canonical
images of the old indices zero and one.  The old-slot equations therefore
preserve the seed length, input coloring, empty stored and connecting-edge
sets, all absent optional values, and the initial closure.  These fields give
the initial constructor for \(A^+\).

Consider an old SE constructor at positive slot \(i.succ\).  The preceding
coloring, preceding tree, current coloring, current data fields, and output
tree are all the corresponding fields of \(A^+\) by the old-slot equations.
The controlled stored-color set at the preceding slot is unchanged by
\(\noderef{TashkinovStoredColorsAppend}\).  Transport, as one dependent
package across these equalities, the maximum-defective choice, the universal
SE alternative, all current-field equations, the exact adjoined tree, its
one-edge-adjoin certificate, its closure, ordered extension, and strict length
increase.  This produces witnesses of exactly the types required by the
lifted SE constructor.  If an RE certificate existed in the appended history,
the appended-to-old map of
\(\noderef{RevisitingExtensionCertificateAppend}\) would give an old RE
certificate, contradicting the old no-RE premise.  Thus the lifted SE
constructor has every required field.

For a PE constructor, use the same equalities and stored-color invariance.
Pointwise in a selected coloring, vertex, and color,
\(\noderef{ParallelExtensionColorAppend}\) identifies the enlarged PE color
rule with the old one; functional and propositional extensionality therefore
identify the two three-argument rule predicates.  Transport in one dependent
package the maximum-defective choice, failure witness for the SE alternative,
the prescribed gamma and its rule, the second stable coloring, its properness
and one-uncoloredness, its stability, the exact exit witness, the component
switch equation, properness and one-uncoloredness after switching, all eight
current fields, the closure, ordered extension, and strict growth.  Because
the rule predicate is transported together with the choice, its occurrences
of the selected coloring and supporting vertex remain coherent.  The same
appended-to-old certificate map proves the lifted no-RE premise.  The
transported package is therefore precisely a PE constructor for \(A^+\).

Finally, suppose the old constructor is RE.  Apply the old-to-appended map in
\(\noderef{RevisitingExtensionCertificateAppend}\) to its available
certificate.  That map preserves gamma, delta, the exact connecting edge,
and both endpoints.  Consequently the old equations for the current coloring,
stored pair, singleton edge set, kind, vertices, and connecting colors rewrite
directly to the equations required for the new certificate.  The old-slot
tree and coloring equations likewise preserve the exact adjoin operation,
closure, ordered extension, and strict growth.  These data form the lifted RE
constructor.  In all four cases its result index is
\((i^+).succ=(i.succ)^+\), or zero in the initial case, which proves the
claim for the original \(j\).
\end{proof}
